\documentclass[../cours_tle.tex]{subfiles}
\begin{document}
\setcounter{chapter}{25}
\chapter{Arguments, exponentielle imaginaire et trigonométrie}
\label{ch:26}
\begin{objectif}
Mathématiques expertes. Choisir la représentation adaptée à un produit ou une puissance. Prérequis : chapitres 17 et 25.
\end{objectif}
\section{Argument et forme trigonométrique}
Pour $z\ne0$, le point $z/|z|$ est sur le cercle unité. Il existe donc $\theta$ tel que
$z=|z|(\cos\theta+i\sin\theta)$.
Un argument est un tel $\theta$ ; tous les arguments diffèrent d'un multiple de $2\pi$. Le nombre $0$ n'a pas d'argument.
Ainsi $1+i=\sqrt2(\cos(\pi/4)+i\sin(\pi/4))$.
Le quotient partie imaginaire/partie réelle ne suffit pas à choisir le quadrant : $1+i$ et $-1-i$ ont le même quotient mais des arguments différents modulo $2\pi$.
\section{Définir l'exponentielle imaginaire}
\begin{definition}
Pour $\theta\in\R$, $e^{i\theta}=\cos\theta+i\sin\theta$.
C'est ici une définition sur l'axe imaginaire ; son lien avec les séries exponentielles sera approfondi en licence.
\end{definition}
Les formules d'addition du chapitre 17 prouvent par multiplication que
$e^{i\alpha}e^{i\beta}=e^{i(\alpha+\beta)}$.
L'ensemble $\mathbb U=\{z:|z|=1\}$ est stable par produit et inverse ; $\overline{e^{i\theta}}=e^{-i\theta}=1/e^{i\theta}$.
Si $z=re^{i\theta}$ et $w=se^{i\varphi}$ sont non nuls,
$zw=rse^{i(\theta+\varphi)}$ et $z/w=(r/s)e^{i(\theta-\varphi)}$.
\section{Euler et Moivre, preuves}
Additionner et soustraire $e^{i\theta}=\cos\theta+i\sin\theta$ et son conjugué donne
\[
\cos\theta=\frac{e^{i\theta}+e^{-i\theta}}2,\qquad
\sin\theta=\frac{e^{i\theta}-e^{-i\theta}}{2i}.
\]
Par récurrence sur la relation multiplicative, puis par inverse pour les entiers négatifs,
$(e^{i\theta})^n=e^{in\theta}$ pour $n\in\Z$.
C'est la formule de Moivre : $(\cos\theta+i\sin\theta)^n=\cos(n\theta)+i\sin(n\theta)$.
\begin{exemple}
$(1+i)^8=(\sqrt2)^8 e^{2i\pi}=16$.
Développer $(\cos x+i\sin x)^3$ et comparer les parties réelles donne
$\cos3x=4\cos^3x-3\cos x$.
Donc $\cos^3x=(3\cos x+\cos3x)/4$, et une primitive est
$3\sin x/4+\sin3x/12$.
\end{exemple}
\section{Linéariser et factoriser}
Les formules d'Euler transforment une puissance trigonométrique en somme d'exponentielles. Par exemple
$\cos^2x=(1+\cos2x)/2$ et
$\sin^2x=(1-\cos2x)/2$.
En additionnant les formules de cosinus pour $(a+b)$ et $(a-b)$,
$\cos a\cos b=(\cos(a+b)+\cos(a-b))/2$.
Ces transformations sont utiles pour intégrer ou étudier une somme d'oscillations.
\section{Un exemple de somme finie}
Si $e^{it}\ne1$, la somme géométrique donne
$\sum_{k=0}^{n-1}e^{ikt}=(1-e^{int})/(1-e^{it})$.
Prendre les parties réelle et imaginaire fournit des sommes de cosinus et de sinus.
Si $t$ est multiple de $2\pi$, la somme vaut directement $n$ : on ne divise pas par zéro pour obtenir ce cas.
\begin{finchapitre}
\section{Utilité et suite du parcours}
Les exponentielles imaginaires sont les briques des séries de Fourier et de l'analyse du signal. Les rotations du plan seront représentées aussi par des matrices.
\begin{retenir}
Les arguments s'additionnent modulo $2\pi$. Pour une somme, préférer souvent la forme algébrique ; pour un produit, quotient ou puissance, la forme exponentielle.
\end{retenir}
\end{finchapitre}
\end{document}
